\documentclass[11pt,a4paper]{article}
\usepackage[margin=24mm]{geometry}
\usepackage{amsmath,amssymb,booktabs,hyperref}
\title{Straight-line residual inspection\\Mathematical and implementation guide, v1}
\author{wenu-imaging}
\date{4 November 2026}
\setlength{\emergencystretch}{2em}
\begin{document}
\maketitle
\noindent This guide describes the completed display algorithm derived from the
owner-approved design \path{docs/developer/design/linearity-residual-display-v1.md}.
It operates on saved ROI means, not individual image pixels. It produces
exploratory diagnostics, not physical shutter or amplifier calibrations.

\section{Inputs and selected population}
For one detector and ROI, each individual flat exposure supplies exposure time
$t_i$ in seconds and saved bias-subtracted mean signal $S_i$ in ADU. Both members
of each pair remain separate. Their common accepted pixel population, source
asset/HDU, acquisition index and original table row remain attached.
There is no overscan correction or additional flat-fielding in this display.
The reference-line population $I$ contains finite $(t_i,S_i)$ and satisfies the
user's inclusive original-signal interval. Blank limits mean all available data.
No weighting, outlier rejection or hidden exposure cut is applied. Points outside
$I$ remain visible and receive predictions if a line is available.

\section{Centered ordinary least squares}
With $n=|I|$, define
\[
 \bar t=\frac1n\sum_{i\in I}t_i,\qquad
 \bar S=\frac1n\sum_{i\in I}S_i,
\]
\[
 b=\frac{\sum_{i\in I}(t_i-\bar t)(S_i-\bar S)}
          {\sum_{i\in I}(t_i-\bar t)^2},\qquad
 a=\bar S-b\bar t.
\]
The slope $b$ has units ADU/s and $a$ has units ADU. A line requires at least two
points and nonzero exposure-time spread. Invalid intervals, insufficient points,
zero time spread and numerical overflow give explicit unavailable states.
The implementation uses compensated sums. Means use the first sample as an
origin, summing scaled differences; predictions use
\[
 \widehat S_i=\bar S+b(t_i-\bar t)
\]
rather than directly evaluating $a+bt_i$, reducing cancellation for large time
origins. No coefficient uncertainties are inferred from these descriptive fits.

\section{Residuals, rates and display}
For each available original measurement, including outside-interval points,
\[
 r_i=S_i-\widehat S_i,\qquad
 r_{\%,i}=100\frac{r_i}{\widehat S_i}\quad(\widehat S_i>0),\qquad
 R_i=\frac{S_i}{t_i}\quad(t_i>0).
\]
Residuals in ADU remain available at nonpositive predictions, but percentage
residuals do not. Zero-time count rates and unavailable signals remain null;
negative measured signals are retained. Nonfinite derived values are null.
Residual plots use original signal or exposure time on the horizontal axis, and
percentage or ADU on the vertical axis, with a zero reference. Curves connect
measurements for inspection; fitting limits affect the line, not displayed data.
The per-ROI summary reports $a,b,n$, the actual selected time/signal ranges and
status. The JSON download retains settings, fits, rows and ASDF source references.
Original ASDF measurements are never rewritten.

\section{What a shutter error can show}
For an otherwise linear, correctly bias-subtracted response under stable light,
\[
 S_i=A(t_i+\delta)\quad\Longrightarrow\quad a=A\delta,\quad b=A,
 \quad R_i=A\left(1+\frac{\delta}{t_i}\right).
\]
A free-intercept line absorbs the constant shutter offset: residuals alone do not
recover it. Its fractional count-rate effect grows as $1/t_i$, motivating the
separate rate view. An additive residual bias also produces a $1/t_i$ term;
lamp variation and detector nonlinearity can contribute structure. No physical
shutter estimate is extracted here. Different signal intervals can yield different
residual baselines even for the same underlying response.

\section{Connection to implementation and validation}
\begin{description}
\item[Arithmetic:] \path{tools/pair-qc-residuals.js}, function
\texttt{QcResiduals.analyze}, taking sample rows and limits, performs grouping, interval
selection, OLS and per-row derived values without modifying its input.
\item[Presentation:] \path{tools/pair-qc-viewer.html} manages metric/axis/interval
controls, ROI visibility, summaries, exact inspection and diagnostic JSON export.
\item[Generation:] \path{tools/view-pair-qc.py} embeds the dependency-free module
into a standalone HTML document and freezes its source with the other inputs.
The JSON alongside the HTML is the saved measurement payload; browser-generated
fit results are separately downloadable.
\item[Checks:] \path{tools/tests/test_pair_qc_residuals.cjs} verifies analytical
solutions, restricted intervals, rank failures, large time origins, null and signed
data, and immutability. \path{tools/tests/test_pair_qc_js.cjs} exercises the actual
embedded script with a minimal DOM surface; it does not verify browser painting.
Existing Python viewer tests check ASDF/geometry/source packaging and preservation.
\end{description}
\end{document}
